\documentclass[11pt]{article}
\usepackage[utf8]{inputenc}
\usepackage[T1]{fontenc}
\usepackage{amsmath,amssymb}
\usepackage{graphicx}
\usepackage{hyperref}
\usepackage{booktabs}
\usepackage{url}
\usepackage{geometry}
\geometry{margin=1in}

\title{Progee v2: An Evidence-First, State-Machine-Governed System for AI-Assisted Software Engineering}
\author{TBD}
\date{January 2026}

\begin{document}
\maketitle
\begin{abstract}
AI-assisted software engineering can reduce implementation time but also introduces new failure modes: unverifiable claims, weak test oracles, and opaque decision-making that makes audits difficult. We present \textbf{Progee v2}, a prototype system that governs AI-assisted development using (1) an explicit task state machine that gates progress through verification steps, and (2) evidence-first persistence of artifacts and checks in a database. The system is implemented in Delphi (FMX) with PostgreSQL storage and DUnitX tests. We release reproducible evaluation artifacts based on end-to-end (E2E) tests that generate machine-readable XML reports, and we discuss current limitations, including an incomplete mutation-testing executor.
\end{abstract}

\begin{quote}
Paper type: System / Tool
Draft date: 2026-01-26
\end{quote}

\section{1. Introduction}
\subsection{1.1 One-sentence positioning (PAPER-001)}
\textbf{Progee v2 is a governance-oriented AI-assisted development system that makes each task’s progress, verification, and failure context auditable by design via an explicit state machine and evidence-first storage.}

\subsection{1.2 Contributions (PAPER-001)}
This paper’s contributions are:
\begin{enumerate}
\item \textbf{Evidence-first artifact persistence for auditability}, including redaction support and bounded evidence snippets for UI/prompt contexts.
\end{enumerate}
\begin{itemize}
\item Implementation: \texttt{CtrlEvidence.pas}, \texttt{sql/16\_context\_evidence\_objects.sql}
\end{itemize}
\begin{enumerate}
\item \textbf{A task state machine that enforces gated verification steps}, including an explicit rework loop with persisted failure history (\texttt{tasks.failure\_context}) that can be summarized back into prompts.
\end{enumerate}
\begin{itemize}
\item Implementation: \texttt{CtrlTaskStateMachine.pas}, \texttt{CtrlWorkshopThread.pas}, \texttt{sql/17\_tasks\_failure\_context.sql}
\end{itemize}
\begin{enumerate}
\item \textbf{A reproducibility package for machine-readable evaluation evidence}, centered on DUnitX XML outputs for E2E tests and a replication guide.
\end{enumerate}
\begin{itemize}
\item Implementation: \texttt{scripts/run\_e2e\_tests.ps1}, \texttt{docs/Replication.md}
\end{itemize}
\begin{enumerate}
\item \textbf{A tested model escalation policy module ("model racing")} that maps rework counts to model levels (policy), intended to be wired into the execution path.
\end{enumerate}
\begin{itemize}
\item Implementation: \texttt{CtrlModelRacing.pas}, \texttt{tests/IntegrationModelRacing.pas}, \texttt{sql/05\_seed\_data.sql}
\end{itemize}
\begin{enumerate}
\item \textbf{A prototype mutation-testing gate integrated into the governance pipeline}, with persisted reports and thresholds, while explicitly acknowledging the executor is not yet production-grade.
\end{enumerate}
\begin{itemize}
\item Implementation: \texttt{sql/11\_mutation\_testing.sql}, \texttt{CtrlMutationEngine.pas}, \texttt{CtrlMutationOperators.pas} (contains TODO/placeholder logic)
\end{itemize}

\subsection{1.3 Contribution-to-evidence mapping (PAPER-004)}
For each contribution, we provide a minimal verification path (commands and artifacts):
\begin{itemize}
\item \textbf{C1 Evidence objects}: run E2E tests (Section 6.2) to populate \texttt{progee\_app\_self.context\_evidence\_objects}, then export \texttt{eval/evidence\_counts.csv} via \texttt{scripts/export\_eval\_metrics.ps1} or query \texttt{scripts/eval\_queries.sql} (section 5/6). Optionally inspect individual rows (type, sha256, redacted).
\item \textbf{C2 State machine + rework failure history}: run E2E tests, then query \texttt{progee\_app\_self.tasks.status}, \texttt{rework\_count}, and \texttt{failure\_context} (see \texttt{scripts/eval\_queries.sql}, section 4). Failure history is summarized into prompts by \texttt{CtrlWorkshopThread.pas}.
\item \textbf{C3 Reproducibility package}: run \texttt{scripts/run\_e2e\_tests.ps1} → \texttt{bin/dunitx-results.xml} (machine‑readable evidence).
\item \textbf{C4 Model racing policy}: run \texttt{tests/IntegrationModelRacing.pas} (DUnitX) and verify default policy keys in \texttt{progee\_core.system\_config} (seeded by \texttt{sql/05\_seed\_data.sql}).
\item \textbf{C5 Mutation-testing gate (prototype)}: run E2E tests, then query \texttt{progee\_app\_self.mutation\_test\_results} / \texttt{mutation\_details}; export summaries via \texttt{scripts/export\_eval\_metrics.ps1} (\texttt{mutation\_score\_summary.csv}, \texttt{mutation\_runs\_by\_status.csv}).
\end{itemize}

\subsection{1.4 Scope}
This paper is a \textbf{system/tool} paper. We emphasize implemented mechanisms and reproducible artifacts over broad causal claims about defect rates.

\section{2. Motivation and Problem Statement}
\subsection{2.1 Why governance, not only generation}
LLM-based code generation can be fast, but correctness and safety depend on:
\begin{itemize}
\item explicit, enforceable verification gates;
\item auditable evidence of what was checked;
\item robust handling of failures and rework.
\end{itemize}

\subsection{2.2 Target user workflow}
Progee v2 targets workflows where:
\begin{itemize}
\item tasks are decomposed into verifiable steps;
\item the system must provide evidence for acceptance decisions;
\item regression is caught by automated tests and (optionally) mutation testing.
\end{itemize}

\section{3. System Overview}
\subsection{3.1 Core entities and persistence}
At a high level, Progee v2 persists:
\begin{itemize}
\item tasks and their states (state machine);
\item artifacts and evidence objects (e.g., test reports, logs, code snippets);
\item model selection decisions;
\item verification outcomes.
\end{itemize}

Key design principle: \textbf{the database is the source of truth}, so the UI and prompts can be derived from persisted evidence rather than ephemeral chat history.

\subsection{3.2 State machine and gates}
The task lifecycle is governed by explicit states (including rework and verification stages).
\begin{itemize}
\item Implementation anchor: \texttt{CtrlTaskStateMachine.pas}
\end{itemize}

\subsection{3.3 Evidence-first collection}
Evidence objects are stored with controlled size (snippets) and optional redaction.
\begin{itemize}
\item Implementation anchor: \texttt{CtrlEvidence.pas}
\end{itemize}

\section{4. Governance Mechanisms}
\subsection{4.1 Rework loop with recorded failure context}
When a task fails a verification gate, Progee records failure context and requests rework.
\begin{itemize}
\item Implementation anchor: \texttt{CtrlTaskStateMachine.pas} (failure\_context JSON)
\end{itemize}

\subsection{4.2 Model escalation policy ("model racing")}
Progee v2 includes a policy module that maps rework counts to model tiers (e.g., level-1 → level-2 at 2 reworks, level-2 → level-3 at 4 reworks). This module is validated by integration tests; wiring it into the AI adapter / per-task execution path is future work.
\begin{itemize}
\item Implementation anchor: \texttt{CtrlModelRacing.pas}, \texttt{tests/IntegrationModelRacing.pas}
\end{itemize}

\subsection{4.3 Verification pipeline (tests; optional mutation gate)}
The workshop pipeline integrates generation and verification steps and persists results.
\begin{itemize}
\item Implementation anchor: \texttt{CtrlWorkshopThread.pas}
\end{itemize}

Mutation testing is represented as a gate with schema and plumbing, but the executor is currently incomplete.
\begin{itemize}
\item Implementation anchor: \texttt{CtrlMutationEngine.pas}, \texttt{CtrlMutationOperators.pas}
\end{itemize}

\section{5. Implementation}
\subsection{5.1 Technology stack}
\begin{itemize}
\item Delphi / FireMonkey (FMX) application
\item PostgreSQL persistence
\item DUnitX for testing
\end{itemize}

\subsection{5.2 Database schema}
Key schema files:
\begin{itemize}
\item Evidence objects: \texttt{sql/16\_context\_evidence\_objects.sql}
\item Mutation testing: \texttt{sql/11\_mutation\_testing.sql}
\end{itemize}

\subsection{5.3 Test harness}
The repo provides a DUnitX test runner:
\begin{itemize}
\item \texttt{bin/ProgeeTests.exe}
\end{itemize}

\section{6. Evaluation (reproducible evidence-first)}
\subsection{6.1 What we evaluate}
We focus on \textbf{reproducibility of verification evidence}:
\begin{itemize}
\item E2E test suite execution
\item Machine-readable XML outputs
\end{itemize}

We avoid claiming defect-rate improvements unless fully supported by reproducible artifacts.

\subsection{6.2 How to reproduce}
See \texttt{docs/Replication.md}. Minimal steps:
\begin{enumerate}
\item Configure a dedicated PostgreSQL database whose name contains \texttt{\_test}.
\item Set environment variables including \texttt{PROGEE\_ALLOW\_DESTRUCTIVE\_TESTS=1}.
\item Run:
\end{enumerate}

\begin{verbatim}
.\scripts\run_e2e_tests.ps1
\end{verbatim}

Expected output:
\begin{itemize}
\item \texttt{bin/dunitx-results.xml}
\end{itemize}

\subsection{6.3 Optional metrics extracted from DB}
For plots/tables, export minimal metrics from DB:
\begin{itemize}
\item number of reworks per task
\item evidence object counts and size distribution
\item mutation score (if available)
\item AI-call volume and token usage (if logs are enabled)
\end{itemize}

We provide:
\begin{itemize}
\item \texttt{scripts/eval\_queries.sql} (ad-hoc queries)
\item \texttt{scripts/export\_eval\_metrics.ps1} → outputs CSVs to \texttt{eval/}
\end{itemize}
\subsection{6.4 Evidence-first end-to-end example (ENG-PAPER-001)}
We provide a minimal, end-to-end evidence chain that can be reproduced from the system:
\begin{enumerate}
\item \textbf{Task → evidence\_ref}: during \texttt{quality\_check} or \texttt{mutation\_test}, Progee persists a validation or mutation report into \texttt{progee\_app\_self.context\_evidence\_objects} and stores its \texttt{id} as \texttt{evidence\_ref} in \texttt{tasks.failure\_context}.
\item \textbf{evidence\_ref → snippet}: the UI’s Evidence tab uses bounded snippet retrieval (\texttt{tail} / \texttt{grep}) via \texttt{CtrlEvidence.pas} to retrieve a safe excerpt of the referenced evidence object.
\item \textbf{snippet → rework prompt}: on rework, \texttt{CtrlWorkshopThread.BuildFailureContextForPrompt} injects the failure summary plus the snippet back into the AI prompt to avoid repeating past mistakes.
\item \textbf{audit record}: each status transition and mutation-gate decision is recorded in \texttt{progee\_app\_self.audit\_logs}.
\end{enumerate}

Minimal reproduction (no manual UI):
\begin{itemize}
\item Run E2E tests (\texttt{scripts/run\_e2e\_tests.ps1}) to populate tasks, failure\_context, and evidence objects.
\item Query \texttt{context\_evidence\_objects} and \texttt{tasks.failure\_context} (see \texttt{scripts/eval\_queries.sql}).
\item Verify audit entries in \texttt{audit\_logs} for task status changes and mutation-gate decisions.
\end{itemize}

\section{7. Limitations \& Threats to Validity (PAPER-005)}
\begin{itemize}
\item \textbf{Mutation testing executor is not production-grade yet}: \texttt{CtrlMutationOperators.pas} contains placeholder/TODO logic for running tests and compile checks; the paper must not claim strong empirical benefits from mutation testing at this stage.
\item \textbf{Model escalation is currently a policy module}: model racing logic is implemented and tested, but the current AI provider/model selection is primarily global-config driven (via \texttt{CtrlAIAdapter.pas}). Wiring escalation into per-task execution is future work.
\item \textbf{Configuration split}: some thresholds and settings are stored in different tables/schemas (e.g., \texttt{progee\_core.system\_config} vs app-level tables). Replication should explicitly document which keys are required.
\item \textbf{External validity}: E2E tests in this repo demonstrate the system on the included workflows; applicability to other languages/toolchains requires adapters.
\item \textbf{Dependence on test quality}: like all testing-based verification, the strength of evidence depends on the quality of written tests.
\end{itemize}

\section{8. Related Work (PAPER-003)}
\textbf{Specifications / contracts.} Formal specification and correctness reasoning date back to Hoare’s axiomatic basis for program correctness, while Design by Contract makes pre/postconditions a first-class discipline for software construction. JML provides a practical specification language and tool ecosystem for Java that operationalizes contract-like design by integrating runtime checking and verification tools.

\textbf{Traceability, auditability, provenance.} Requirements traceability has long been recognized as a core engineering problem, with classic analyses distinguishing pre‑RS and post‑RS traceability and modern books consolidating practice across the lifecycle. Recent work on traceability in the wild highlights real-world gaps in link completeness. For provenance and supply-chain auditability, the W3C PROV family formalizes provenance representations, while SLSA provides a practical, staged framework for build provenance and verification.

\textbf{Mutation testing.} Mutation testing has an extensive body of theory, tools, and empirical studies, including major surveys and practical systems (e.g., µJava) that operationalize mutation operators and scoring for test adequacy.

\textbf{AI-assisted software engineering \& agentic verification.} Code LMs (Codex, AlphaCode) show strong synthesis ability but still require robust evaluation and governance. SWE-bench and SWE-agent provide repository-scale benchmarks and agent interfaces. A growing body of agentic and self-critique methods (ReAct, Reflexion, Self‑Refine) explores how models can plan, act, and refine; alignment‑oriented ideas like Constitutional AI and Debate provide mechanisms for structured critique or oversight. Progee v2 complements these lines with an explicit, auditable state machine and evidence‑first persistence.

\textbf{References (selected)}  
\begin{enumerate}
\item C. A. R. Hoare. \textit{An Axiomatic Basis for Computer Programming}. Communications of the ACM, 1969.  
\item B. Meyer. \textit{Applying Design by Contract}. IEEE Computer, 25(10), 1992.  
\item L. Burdy et al. \textit{An Overview of JML Tools and Applications}. Software Tools for Technology Transfer, 7(3), 2005.  
\item O. Gotel \& A. Finkelstein. \textit{An Analysis of the Requirements Traceability Problem}. ICRE/RE, 1994.  
\item J. Cleland-Huang, O. Gotel, A. Zisman (eds.). \textit{Software and Systems Traceability}. Springer, 2012.  
\item M. Rath et al. \textit{Traceability in the Wild: Automatically Augmenting Incomplete Trace Links}. 2018.  
\item W3C. \textit{PROV‑N: The Provenance Notation}. W3C Recommendation, 2013.  
\item SLSA Community. \textit{SLSA v1.1 Specification} (Approved). 2025.  
\item B. Kitchenham, T. Dybå, M. Jørgensen. \textit{Evidence‑Based Software Engineering}. ICSE, 2004.  
\item Y. Jia \& M. Harman. \textit{An Analysis and Survey of the Development of Mutation Testing}. IEEE TSE, 37(5), 2011.  
\item M. Papadakis et al. \textit{Mutation Testing Advances: An Analysis and Survey}. Advances in Computers, 2019.  
\item Y.‑S. Ma, J. Offutt, Y.‑R. Kwon. \textit{µJava: An Automated Class Mutation System}. STVR, 15(2), 2005.  
\item M. Chen et al. \textit{Evaluating Large Language Models Trained on Code (Codex)}. arXiv:2107.03374, 2021.  
\item Y. Li et al. \textit{Competition‑Level Code Generation with AlphaCode}. arXiv:2203.07814, 2022.  
\item C. E. Jimenez et al. \textit{SWE‑bench: Can Language Models Resolve Real‑World GitHub Issues?} ICLR, 2024.  
\item J. Yang et al. \textit{SWE‑agent: Agent‑Computer Interfaces Enable Automated Software Engineering}. arXiv:2405.15793, 2024.  
\item S. Yao et al. \textit{ReAct: Synergizing Reasoning and Acting in Language Models}. arXiv:2210.03629, 2022.  
\item N. Shinn et al. \textit{Reflexion: Language Agents with Verbal Reinforcement Learning}. arXiv:2303.11366, 2023.  
\item A. Madaan et al. \textit{Self‑Refine: Iterative Refinement with Self‑Feedback}. arXiv:2303.17651, 2023.  
\item G. Irving, P. Christiano, D. Amodei. \textit{AI Safety via Debate}. arXiv:1805.00899, 2018.  
\item Y. Bai et al. \textit{Constitutional AI: Harmlessness from AI Feedback}. arXiv:2212.08073, 2022.  
\end{enumerate}

\section{9. Conclusion and Future Work}
We presented Progee v2 as an evidence-first, state-machine-governed system for AI-assisted software engineering, and provided a reproducible evaluation path via E2E tests with XML outputs.

Future work:
\begin{itemize}
\item complete and validate a minimal mutation testing executor for Delphi/DUnitX workflows
\item add standardized DB export scripts for evaluation metrics
\end{itemize}

\section{Artifact Availability}
\begin{itemize}
\item Replication guide: \texttt{docs/Replication.md}
\item One-command E2E runner: \texttt{scripts/run\_e2e\_tests.ps1}
\item DUnitX XML output (generated): \texttt{bin/dunitx-results.xml}
\end{itemize}

\end{document}
